\chapter{The Learning Problem}
\lecture{1}{9 Sep.}{}

\section{Course Introduction}

\source{slides p.1--5}

Machine learning is a mixture of theoretical and practical tools, so any course
about it has to decide how to mix them. There are two natural extremes, and both
of them fail in their own way.

\begin{itemize}
    \item \emph{Theory oriented}: derive everything deeply, so that every
        statement rests on a proof and the understanding is solid. The cost is
        that it is less interesting to a general audience --- one may finish such
        a course able to prove a generalization bound, yet unable to fit a model
        to a real data set.
    \item \emph{Techniques oriented}: flash over the sexiest techniques broadly,
        so that the coverage looks shiny. The cost is that there are simply too
        many techniques: they become hard to choose between, and hard to use
        properly, because nothing tells you \emph{why} one of them should work
        here and not there.
\end{itemize}

\begin{moral}
    Our approach: \vocab{foundation oriented}.
\end{moral}

By \emph{foundation} we mean the part that does not change when the fashionable
technique changes. Concretely, the course is a mixture of philosophical
illustrations, key theory, core techniques, usage in practice, and hopefully
jokes; the selecting principle is always \emph{what every machine learning user
should know}. The payoff of insisting on foundations is that they transfer: a
student who owns them can pick up a future or untaught technique, or go read
deeper theory, without being taught each one explicitly.

The material is organized \vocab{story-like}, as a sequence of questions, each
one asked only after the previous one has been answered:

\begin{enumerate}
    \item \textbf{When} can machines learn? (illustrative + technical)
    \item \textbf{Why} can machines learn? (theoretical + illustrative)
    \item \textbf{How} can machines learn? (technical + practical)
    \item \textbf{How} can machines learn \textbf{better}? (practical + theoretical)
\end{enumerate}

\begin{note}
    The four questions above are the backbone of the \emph{foundations} half of
    the story. The NTU version continues it with three more chapters, each one a
    different answer to ``where do good features come from'':
    \begin{enumerate}[start=5]
        \item How can machines learn by \emph{embedding numerous} features?
            (support vector machines)
        \item How can machines learn by \emph{combining predictive} features?
            (aggregation, boosting, trees and forests)
        \item How can machines learn by \emph{distilling hidden} features?
            (neural networks and deep learning)
    \end{enumerate}
\end{note}

This first chapter sits entirely inside question~1, and it is the most basic
piece of it: before asking when machines \emph{can} learn, we must say what the
learning problem even \emph{is}. That is the plan of the chapter: we first ask
what machine learning is (\autoref{sec:what-is-ml}), then look at where it
already lives (\autoref{sec:applications}), then cut it into its formal
components (\autoref{sec:components}), and finally place it next to the
neighbouring fields it is often confused with (\autoref{sec:other-fields}).

\begin{exercise}[Fun Time]
    Which of the following descriptions of this course is true?
    \begin{enumerate}[(1)]
        \item the course will be taught in Taiwanese;
        \item the course will tell me the techniques that create the android
            Lieutenant Commander Data in \emph{Star Trek};
        \item the course will be 15 weeks long;
        \item the course will be story-like.
    \end{enumerate}
\end{exercise}

\begin{proof}[Reference Answer]
    (4). As for the others: (1) no, the instructor's Taiwanese is unfortunately
    not good enough for teaching (yet); (2) no, although what we teach may serve
    as foundations of those (future) techniques; (3) no, unless you choose to
    join the next course. So: let us begin the story.
\end{proof}

\section{What is Machine Learning}\label{sec:what-is-ml}

\source{slides p.6--11}

\subsection{From Learning to Machine Learning}

The safest way into a definition is by analogy with the thing we already do
ourselves.

\begin{definition}[Learning]
    \vocab{Learning} is acquiring \emph{skill} with experience accumulated from
    \emph{observations}.
\end{definition}

\begin{center}
    \begin{tikzpicture}[node distance=1.3cm]
        \node[mlbox] (a) {observations};
        \node[mlop, right=of a] (b) {learning};
        \node[mlbox, right=of b] (c) {skill};
        \draw[mlarrow] (a) -- (b);
        \draw[mlarrow] (b) -- (c);
    \end{tikzpicture}
\end{center}

A machine has no eyes and no childhood, so the only observations it can
accumulate are the ones we hand it in recorded form. Replacing
\emph{observations} by \emph{data} and \emph{learning} by \emph{computation}
gives the machine version.

\begin{definition}[Machine learning, first attempt]
    \vocab{Machine learning} is acquiring skill with experience
    accumulated/computed from \emph{data}.
\end{definition}

\begin{center}
    \begin{tikzpicture}[node distance=1.3cm]
        \node[mlbox] (a) {data};
        \node[mlop, right=of a] (b) {ML};
        \node[mlbox, right=of b] (c) {skill};
        \draw[mlarrow] (a) -- (b);
        \draw[mlarrow] (b) -- (c);
    \end{tikzpicture}
\end{center}

The word \emph{data} is concrete enough to be useful. The word \emph{skill} is
not: it is exactly the kind of word that cannot be handed to a computer, so the
definition is not finished yet.

\begin{question}
    What is skill?
\end{question}

\subsection{A More Concrete Definition}

\begin{intuition}
    skill \(\iff\) improve some \emph{performance measure}
    (e.g.\ prediction accuracy).
\end{intuition}

This is the move that makes the whole subject possible: a vague goal (``be
good at recognizing trees'') is replaced by a number that can be computed and
watched (``fraction of tree images classified correctly''). Every later chapter
is, in one way or another, an argument about which number to watch and what
watching it really guarantees.

\begin{definition}[Machine learning]\label{def:ml-performance}
    \vocab{Machine learning} is improving some \emph{performance measure} with
    experience computed from \emph{data}.
\end{definition}

\begin{center}
    \begin{tikzpicture}[node distance=1.3cm]
        \node[mlbox] (a) {data};
        \node[mlop, right=of a] (b) {ML};
        \node[mlbox, right=of b] (c) {improved\\performance\\measure};
        \draw[mlarrow] (a) -- (b);
        \draw[mlarrow] (b) -- (c);
    \end{tikzpicture}
\end{center}

\begin{eg}[An application in computational finance]
    Take the data to be historical stock data, and the performance measure to be
    investment gain. Then machine learning here means: use the stock data to
    compute a trading rule whose gain is higher than what we had before.
    \begin{center}
        \begin{tikzpicture}[node distance=1.3cm]
            \node[mlbox] (a) {stock data};
            \node[mlop, right=of a] (b) {ML};
            \node[mlbox, right=of b] (c) {more\\investment gain};
            \draw[mlarrow] (a) -- (b);
            \draw[mlarrow] (b) -- (c);
        \end{tikzpicture}
    \end{center}
\end{eg}

\subsection{Why Use Machine Learning?}

We now know what machine learning \emph{is}. A separate matter is why one would
ever want it, given that a computer can already be programmed directly.

\begin{question}
    Why use machine learning?
\end{question}

\begin{eg}[Tree recognition]
    Suppose we want a system that looks at a photograph and answers whether
    there is a tree in it.
    \begin{itemize}
        \item To \emph{`define' trees and hand-program} the answer is difficult:
            try writing down, in code, a condition that accepts every pine,
            palm and bonsai and rejects every bush and lamppost.
        \item To \emph{learn from data (observations) and recognize} is easy
            enough that a 3-year-old can do it, having never received a
            definition of ``tree'' in their life --- only examples.
    \end{itemize}
    So an `ML-based tree recognition system' can be easier to \emph{build} than
    a hand-programmed one, even though it is harder to \emph{explain}.
\end{eg}

\begin{moral}
    ML is an alternative route to build complicated systems.
\end{moral}

The useful way to read that moral is as a list of situations where the direct
route is blocked:

\begin{itemize}
    \item when humans cannot program the system manually --- navigating on Mars;
    \item when humans cannot `define the solution' easily ---
        speech/visual recognition;
    \item when rapid decisions are needed that humans cannot make ---
        high-frequency trading;
    \item when the system must be user-oriented at a massive scale ---
        consumer-targeted marketing.
\end{itemize}

\begin{remark}
    Give a computer a fish, you feed it for a day; teach it how to fish, you
    feed it for a lifetime. :-)
\end{remark}

\subsection{Key Essence of Machine Learning}

Machine learning is not always the right tool, and the definition above already
tells us when it is not. Reading \autoref{def:ml-performance} backwards, three
things must hold before the phrase even makes sense.

\begin{enumerate}
    \item There exists some \vocab{underlying pattern} to be learned --- so that
        a `performance measure' \emph{can} be improved. Without a pattern there
        is nothing to be better at.
    \item But there is no programmable (easy) \emph{definition} of that pattern
        --- so that `ML' is \emph{needed}. If you can just write the rule down,
        write it down.
    \item Somehow there is \vocab{data} about the pattern --- so that ML has some
        `inputs' to learn from. A pattern nobody has recorded cannot be computed
        from.
\end{enumerate}

\begin{moral}
    Key essence: it helps decide whether to use ML at all.
\end{moral}

\begin{exercise}[Fun Time]
    Which of the following is best suited for machine learning?
    \begin{enumerate}[(1)]
        \item predicting whether the next cry of a baby girl happens at an
            even-numbered minute or not;
        \item determining whether a given graph contains a cycle;
        \item deciding whether to approve a credit card to some customer;
        \item guessing whether the earth will be destroyed by the misuse of
            nuclear power in the next ten years.
    \end{enumerate}
\end{exercise}

\begin{proof}[Reference Answer]
    (3), and the three essences say exactly why each of the others fails.
    \begin{enumerate}[(1)]
        \item Fails essence~1: there is no pattern.
        \item Fails essence~2: the definition is programmable (run a cycle
            detection algorithm).
        \item Passes all three --- pattern: customer behaviour; definition: not
            easily programmable; data: the history of bank operations.
        \item Fails essence~3: there is arguably no (or not enough) data yet.
    \end{enumerate}
\end{proof}

\section{Applications of Machine Learning}\label{sec:applications}

\source{slides p.12--16}

Every application below is the same picture,
\begin{center}
    \begin{tikzpicture}[node distance=1.3cm]
        \node[mlbox] (a) {data};
        \node[mlop, right=of a] (b) {ML};
        \node[mlbox, right=of b] (c) {skill};
        \draw[mlarrow] (a) -- (b);
        \draw[mlarrow] (b) -- (c);
    \end{tikzpicture}
\end{center}
with the two ends filled in differently. So the exercise in reading an
application is always: \emph{what is the data, and what is the skill?}

\subsection{Daily Needs: Food, Clothing, Housing, Transportation}

\begin{enumerate}
    \item \textbf{Food} (Sadilek et al., 2013).
        \emph{Data}: Twitter data (words \(+\) location).
        \emph{Skill}: tell the food-poisoning likeliness of a restaurant
        properly.
    \item \textbf{Clothing} (Abu-Mostafa, 2012).
        \emph{Data}: sales figures \(+\) client surveys.
        \emph{Skill}: give good fashion recommendations to clients.
    \item \textbf{Housing} (Tsanas and Xifara, 2012).
        \emph{Data}: characteristics of buildings and their energy load.
        \emph{Skill}: predict the energy load of other buildings closely.
    \item \textbf{Transportation} (Stallkamp et al., 2012).
        \emph{Data}: some traffic sign images and their meanings.
        \emph{Skill}: recognize traffic signs accurately.
\end{enumerate}

\begin{moral}
    ML is everywhere!
\end{moral}

\subsection{Education}

\emph{Data}: students' records on quizzes in a mathematics tutoring system.
\emph{Skill}: predict whether a student can give a correct answer to
\emph{another} quiz question.

What makes this example instructive is that the pattern one guesses is not a
formula in the recorded variables at all; it is a formula in quantities nobody
ever measured:
\[
    \text{answer correctly} \approx
    \llbracket\, \text{recent strength of student}
    > \text{difficulty of question} \,\rrbracket.
\]

\begin{notation}
    \(\llbracket \cdot \rrbracket\) is the \emph{Iverson bracket}: it equals
    \(1\) when the statement inside is true and \(0\) when it is false.
\end{notation}

Neither ``strength'' nor ``difficulty'' appears in the data --- only who
answered what, and whether it was right. Feeding roughly \(9\) million records
from \(3000\) students to a learning algorithm lets it
\emph{reverse-engineer} both numbers automatically, purely from the requirement
that the formula above should reproduce the observed answers.

\begin{moral}
    This was the key part of the world-champion system from National Taiwan
    University in KDDCup 2010.
\end{moral}

\subsection{Entertainment: Recommender Systems}

\emph{Data}: how many users have rated some movies.
\emph{Skill}: predict how a user would rate an unrated movie.

\begin{eg}[A hot problem]
    \begin{itemize}
        \item A competition held by Netflix in 2006: \(100{,}480{,}507\) ratings
            that \(480{,}189\) users gave to \(17{,}770\) movies, with a
            \(10\%\) improvement worth a one-million-dollar prize.
        \item A similar competition (movies \(\to\) songs) held by Yahoo!\ in
            KDDCup 2011: \(252{,}800{,}275\) ratings that \(1{,}000{,}990\)
            users gave to \(624{,}961\) songs.
    \end{itemize}
\end{eg}

\begin{question}
    How can machines learn our preferences?
\end{question}

The same trick as in the education example applies, and it is worth seeing that
it is literally the same trick. Postulate that a viewer and a movie are each
described by the \emph{same} list of hidden \vocab{factors} --- how much comedy,
how much action, how much of a blockbuster --- read on the viewer's side as
\emph{how much this viewer wants it} and on the movie's side as \emph{how much
this movie has it}.

\begin{center}
    \begin{adjustbox}{max width=\textwidth}
        \begin{tikzpicture}[font=\small,
            slot/.style={draw=black!65, minimum width=0.8cm,
                    minimum height=0.62cm, inner sep=0pt},
            flab/.style={font=\scriptsize, inner sep=1pt},
            tick/.style={-{Stealth[length=1.4mm]}, thin, draw=black!55}]

            % ── the two factor vectors, as rows of slots in matching order
            \foreach \y/\name in {1.5/viewer, -1.5/movie}{
                \node[font=\small, anchor=east] at (-0.15,\y) {\name};
                \foreach \i in {0,1,2}{\node[slot] at ({0.4+0.8*\i},\y) {};}
                \node[slot, minimum width=2.2cm] at (3.5,\y) {\(\cdots\)};
                \node[slot] at (5.0,\y) {};
            }
            % how much the viewer wants each factor (dot size = magnitude)
            \fill[ForestGreen] (0.4, 1.5) circle (0.09);
            \fill[NavyBlue]    (1.2, 1.5) circle (0.17);
            \fill[Magenta]     (2.0, 1.5) circle (0.09);
            \fill[Violet]      (5.0, 1.5) circle (0.07);
            % how much the movie has of each factor
            \fill[ForestGreen] (0.4,-1.5) circle (0.18);
            \fill[NavyBlue]    (1.2,-1.5) circle (0.08);
            \fill[Magenta]     (2.0,-1.5) circle (0.09);
            \fill[Violet]      (5.0,-1.5) circle (0.18);

            % ── what each slot stands for
            \foreach \i/\up/\dn in {%
                0/likes comedy?/comedy content,
                1/likes action?/action content,
                2/prefers blockbusters?/blockbuster?}{
                \node[flab, rotate=58, anchor=south west]
                    at ({0.55+0.8*\i},{1.98+0.2*\i}) {\up};
                \draw[tick] ({0.5+0.8*\i},{1.95+0.2*\i})
                    -- ({0.4+0.8*\i},1.84);
                \node[flab, rotate=58, anchor=north east]
                    at ({0.55+0.8*\i},{-1.98-0.2*\i}) {\dn};
                \draw[tick] ({0.5+0.8*\i},{-1.95-0.2*\i})
                    -- ({0.4+0.8*\i},-1.84);
            }
            \node[flab, rotate=58, anchor=south west]
                at (5.15,1.98) {likes Tom Cruise?};
            \draw[tick] (5.1,1.95) -- (5.0,1.84);
            \node[flab, rotate=58, anchor=north east]
                at (5.15,-1.98) {Tom Cruise in it?};
            \draw[tick] (5.1,-1.95) -- (5.0,-1.84);

            % ── slot k of one row is paired with slot k of the other
            \foreach \x in {1.2, 2.0, 5.0}{
                \draw[black!25, densely dotted] (\x,1.19) -- (\x,-1.19);}
            \draw[{Stealth[length=2mm]}-{Stealth[length=2mm]}, thick, black!70]
                (0.4,1.19) -- (0.4,-1.19);
            \node[font=\scriptsize, align=center, fill=white, inner sep=2pt]
                at (1.7,0) {match movie and\\viewer factors};

            % ── then sum the matched contributions
            \draw[mlarrow] (3.35,0) -- (4.95,0);
            \node[font=\scriptsize, align=center, anchor=south]
                at (4.15,0.16) {add contributions\\from each factor};
            \node[align=center, font=\small\bfseries, color=RawSienna,
                anchor=west] at (5.15,0) {predicted\\rating};
        \end{tikzpicture}
    \end{adjustbox}
\end{center}

Both rows carry the same factors in the same order, so a viewer and a movie
live in one common space; the size of a dot is the size of a number. The
vertical arrow pairs slot \(k\) with slot \(k\), the horizontal one sums the
pairs, and together they are a single formula: writing \(v_k\) for how much the
viewer wants factor \(k\) and \(m_k\) for how much the movie has of it,
\[
    \text{predicted rating} \approx \sum_{k=1}^{K} v_k m_k ,
\]
the inner product of the two factor vectors. The product is the point --- a
factor contributes only where the two sides \emph{agree}, so our
action-loving viewer would rate this comedy rather low.

The dashes say the list is long, and the last slot says how arbitrary it may
be. Nobody drew that list up: as with the students' strength and difficulty,
the factors are never observed, and the learning is what invents them.
\begin{align*}
    \text{pattern:} \quad & \text{rating} \longleftarrow
        \text{viewer/movie factors};\\
    \text{learning:} \quad & \text{known ratings} \longrightarrow
        \text{learned factors} \longrightarrow
        \text{unknown rating prediction}.
\end{align*}

\begin{moral}
    Again the key part of the world-champion system from National Taiwan
    University, this time in KDDCup 2011.
\end{moral}

\begin{exercise}[Fun Time]
    Which of the following fields \emph{cannot} use machine learning?
    \begin{enumerate}[(1)]
        \item Finance;
        \item Medicine;
        \item Law;
        \item none of the above.
    \end{enumerate}
\end{exercise}

\begin{proof}[Reference Answer]
    (4). Finance can predict stock prices from data, medicine can predict
    medicine effects from data, and law can summarize legal documents from data.
    Welcome to study this hot topic.
\end{proof}

\section{Components of Machine Learning}\label{sec:components}

\source{slides p.17--22}

So far ``machine learning'' has been a box in a picture. This section opens the
box and names its parts. The running metaphor is \emph{credit approval}: a bank
receives an application and must decide whether issuing a card to this person is
good for the bank.

\begin{center}
    \begin{tabular}{@{}lr@{}}
        \toprule
        \multicolumn{2}{@{}l}{\textbf{Applicant Information}}\\
        \midrule
        age               & 23 years\\
        gender            & female\\
        annual salary     & NTD 1,000,000\\
        year in residence & 1 year\\
        year in job       & 0.5 year\\
        current debt      & 200,000\\
        \bottomrule
    \end{tabular}
\end{center}

The unknown pattern to be learned is: \emph{is approving this credit card good
for the bank?}

\subsection{Formalizing the Learning Problem}

\begin{notation}
    \hfill
    \begin{itemize}
        \item \textbf{Input}: \(x \in \mathcal{X}\) (a customer application).
        \item \textbf{Output}: \(y \in \mathcal{Y}\) (good/bad after approving
            the credit card).
        \item The unknown pattern to be learned \(\iff\) the
            \vocab{target function}
            \[
                f\colon \mathcal{X} \to \mathcal{Y}
            \]
            (the ideal credit approval formula).
        \item The data \(\iff\) the \vocab{training examples}
            \[
                \mathcal{D} = \left\{ (x_1, y_1), (x_2, y_2), \dots ,
                (x_N, y_N) \right\}
            \]
            (historical records in the bank).
        \item The skill, hopefully with good performance, \(\iff\) the
            \vocab{hypothesis}
            \[
                g\colon \mathcal{X} \to \mathcal{Y}
            \]
            (the `learned' formula to be used).
    \end{itemize}
\end{notation}

\begin{center}
    \begin{tikzpicture}[node distance=1.3cm]
        \node[mlbox, minimum width=3.4cm] (a) {\(\left\{ (x_n, y_n) \right\}\)
            from \(f\)};
        \node[mlop, right=of a] (b) {ML};
        \node[mlbox, right=of b] (c) {\(g\)};
        \draw[mlarrow] (a) -- (b);
        \draw[mlarrow] (b) -- (c);
    \end{tikzpicture}
\end{center}

\begin{remark}
    The asymmetry between \(f\) and \(g\) is the whole difficulty of the
    subject, so it is worth stating plainly.
    \begin{itemize}
        \item \(f\) is \emph{unknown} --- indeed, by essence~2 of the previous
            section, it has no programmable definition. We never see \(f\); we
            only ever see \(N\) of its values.
        \item \(g\) is what we produce, and we hope \(g \approx f\). But \(g\)
            may differ from \(f\), and perfection is `impossible' precisely
            because \(f\) is unknown --- there is no way to check the answer.
    \end{itemize}
    Everything in \emph{Why can machines learn?} is an attempt to say something
    rigorous about \(g \approx f\) despite never seeing \(f\).
\end{remark}

Drawing the arrows in the order in which things actually happen gives the
learning flow.

\begin{center}
    \begin{adjustbox}{max width=\textwidth}
        \begin{tikzpicture}
            \node[mlbox, minimum width=4.2cm] (f) at (0,2.25)
                {unknown target function\\
                 \(f\colon \mathcal{X} \to \mathcal{Y}\)\\[2pt]
                 {\scriptsize (ideal credit approval formula)}};
            \node[mlbox, minimum width=5.2cm] (D) at (0,0)
                {training examples\\
                 \(\mathcal{D}\colon (x_1, y_1), \dots , (x_N, y_N)\)\\[2pt]
                 {\scriptsize (historical records in bank)}};
            \node[mlop, minimum width=2.4cm] (A) at (5.4,0)
                {learning\\algorithm\\\(\mathcal{A}\)};
            \node[mlbox, minimum width=4.4cm] (g) at (10.2,0)
                {final hypothesis\\
                 \(g \approx f\)\\[2pt]
                 {\scriptsize (`learned' formula to be used)}};
            \draw[mlarrow] (f) -- (D);
            \draw[mlarrow] (D) -- (A);
            \draw[mlarrow] (A) -- (g);
        \end{tikzpicture}
    \end{adjustbox}
\end{center}

\subsection{The Learning Model}

The diagram still hides one thing.

\begin{question}
    What does \(g\) look like?
\end{question}

An algorithm cannot search over \emph{all} functions \(\mathcal{X} \to \mathcal{Y}\), so we
must first decide which formulas are even candidates.

\begin{definition}[Hypothesis set]
    The \vocab{hypothesis set} \(\mathcal{H} = \left\{ h_k \right\}\) is the set
    of candidate formulas that the algorithm is allowed to return, so that
    \(g \in \mathcal{H}\).
\end{definition}

\begin{eg}
    For credit approval, one might approve if
    \begin{itemize}
        \item \(h_1\): annual salary \(>\) NTD 800,000;
        \item \(h_2\): debt \(>\) NTD 100,000 \quad (really?);
        \item \(h_3\): year in job \(\le 2\) \quad (really?).
    \end{itemize}
    Note that \(\mathcal{H}\) happily contains bad hypotheses next to good ones.
    That is not a defect: it is up to \(\mathcal{A}\) to pick the `best' one out
    of \(\mathcal{H}\) and call it \(g\).
\end{eg}

\begin{definition}[Learning model]
    A \vocab{learning model} is the pair consisting of the hypothesis set
    \(\mathcal{H}\) and the learning algorithm \(\mathcal{A}\).
\end{definition}

\begin{center}
    \begin{adjustbox}{max width=\textwidth}
        \begin{tikzpicture}
            \node[mlbox, minimum width=4.2cm] (f) at (0,2.25)
                {unknown target function\\
                 \(f\colon \mathcal{X} \to \mathcal{Y}\)\\[2pt]
                 {\scriptsize (ideal credit approval formula)}};
            \node[mlbox, minimum width=5.2cm] (D) at (0,0)
                {training examples\\
                 \(\mathcal{D}\colon (x_1, y_1), \dots , (x_N, y_N)\)\\[2pt]
                 {\scriptsize (historical records in bank)}};
            \node[mlop, minimum width=2.4cm] (A) at (5.4,0)
                {learning\\algorithm\\\(\mathcal{A}\)};
            \node[mlbox, minimum width=4.4cm] (g) at (10.2,0)
                {final hypothesis\\
                 \(g \approx f\)\\[2pt]
                 {\scriptsize (`learned' formula to be used)}};
            \node[mlbox, minimum width=3.4cm] (H) at (5.4,-2.25)
                {hypothesis set\\
                 \(\mathcal{H}\)\\[2pt]
                 {\scriptsize (set of candidate formulas)}};
            \draw[mlarrow] (f) -- (D);
            \draw[mlarrow] (D) -- (A);
            \draw[mlarrow] (A) -- (g);
            \draw[mlarrow] (H) -- (A);
        \end{tikzpicture}
    \end{adjustbox}
\end{center}

This is the picture the rest of the course elaborates. The dashed halves of the
story are already visible in it: \(f\) and \(\mathcal{D}\) are given by the
world and we cannot choose them, while \(\mathcal{H}\) and \(\mathcal{A}\) are
ours to design.

\begin{moral}
    Machine learning: use data to compute a hypothesis \(g\) that approximates
    the target \(f\).
\end{moral}

\begin{exercise}[Fun Time]
    How should the four sets below be used to form a learning problem for song
    recommendation?
    \begin{align*}
        S_1 &= [0, 100];\\
        S_2 &= \text{all possible (userid, songid) pairs};\\
        S_3 &= \text{all formulas that `multiply' user factors and song
            factors,}\\
            &\qquad \text{indexed by all possible combinations of such
            factors};\\
        S_4 &= \text{1,000,000 pairs of ((userid, songid), rating)}.
    \end{align*}
    \begin{enumerate}[(1)]
        \item \(S_1 = \mathcal{X},\ S_2 = \mathcal{Y},\ S_3 = \mathcal{H},\
            S_4 = \mathcal{D}\);
        \item \(S_1 = \mathcal{Y},\ S_2 = \mathcal{X},\ S_3 = \mathcal{H},\
            S_4 = \mathcal{D}\);
        \item \(S_1 = \mathcal{D},\ S_2 = \mathcal{H},\ S_3 = \mathcal{Y},\
            S_4 = \mathcal{X}\);
        \item \(S_1 = \mathcal{X},\ S_2 = \mathcal{D},\ S_3 = \mathcal{Y},\
            S_4 = \mathcal{H}\).
    \end{enumerate}
\end{exercise}

\begin{proof}[Reference Answer]
    (2). Reading the whole problem in one line,
    \[
        S_4 \xrightarrow{\ \mathcal{A} \text{ on } S_3\ }
        \left( g\colon S_2 \to S_1 \right),
    \]
    i.e.\ the input space is the set of (user, song) pairs, the output space is
    the rating scale \([0,100]\), the hypothesis set is the family of
    factor-multiplying formulas, and the data is the list of observed ratings.
\end{proof}

\section{Machine Learning and Other Fields}\label{sec:other-fields}

\source{slides p.23--27}

Machine learning did not grow up alone, and much of the confusion around the
name comes from its neighbours. Comparing the one-line definitions is enough to
see both the overlap and the difference in \emph{focus} --- which is usually
where the real difference lies.

\subsection{Machine Learning and Data Mining}

\begin{center}
    \begin{tabular}{@{}p{0.45\textwidth}p{0.45\textwidth}@{}}
        \toprule
        \textbf{Machine Learning} & \textbf{Data Mining}\\
        \midrule
        use data to compute a hypothesis \(g\) that approximates a target \(f\)
        & use (huge) data to find a property that is \emph{interesting}\\
        \bottomrule
    \end{tabular}
\end{center}

\begin{itemize}
    \item If the `interesting property' \emph{is} the `hypothesis that
        approximates the target', then ML \(=\) DM. This is usually what KDDCup
        does.
    \item If the `interesting property' is merely \emph{related} to such a
        hypothesis, then DM can help ML and vice versa --- often, but not
        always.
    \item Traditional DM also puts weight on efficient computation in large
        databases, which ML traditionally does not.
\end{itemize}

\begin{moral}
    It is difficult to distinguish ML and DM in reality.
\end{moral}

\subsection{Machine Learning and Artificial Intelligence}

\begin{center}
    \begin{tabular}{@{}p{0.45\textwidth}p{0.45\textwidth}@{}}
        \toprule
        \textbf{Machine Learning} & \textbf{Artificial Intelligence}\\
        \midrule
        use data to compute a hypothesis \(g\) that approximates a target \(f\)
        & compute something that shows \emph{intelligent behaviour}\\
        \bottomrule
    \end{tabular}
\end{center}

\begin{itemize}
    \item \(g \approx f\) is something that shows intelligent behaviour, so ML
        can realize AI --- among other routes.
    \item For example, in chess playing: traditional AI builds a game tree,
        whereas ML for AI learns from board data.
\end{itemize}

\begin{moral}
    ML is one possible route to realize AI.
\end{moral}

\subsection{Machine Learning and Statistics}

\begin{center}
    \begin{tabular}{@{}p{0.45\textwidth}p{0.45\textwidth}@{}}
        \toprule
        \textbf{Machine Learning} & \textbf{Statistics}\\
        \midrule
        use data to compute a hypothesis \(g\) that approximates a target \(f\)
        & use data to make \emph{inference} about an unknown process\\
        \bottomrule
    \end{tabular}
\end{center}

\begin{itemize}
    \item \(g\) is an inference outcome and \(f\) is something unknown, so
        statistics can be used to achieve ML.
    \item Traditional statistics also focuses on provable results under
        mathematical assumptions, and cares less about computation.
\end{itemize}

\begin{moral}
    Statistics: many useful tools for ML.
\end{moral}

\begin{exercise}[Fun Time]
    Which of the following claims is \emph{not} totally true?
    \begin{enumerate}[(1)]
        \item machine learning is a route to realize artificial intelligence;
        \item machine learning, data mining and statistics all need data;
        \item data mining is just another name for machine learning;
        \item statistics can be used for data mining.
    \end{enumerate}
\end{exercise}

\begin{proof}[Reference Answer]
    (3). While data mining and machine learning do share a huge overlap, they
    are arguably not equivalent, because of the difference of focus.
\end{proof}

\hrulebar

\subsection*{Summary}

\begin{description}
    \item[Course Introduction] foundation oriented and story-like.
    \item[What is Machine Learning] use data to approximate a target.
    \item[Applications of Machine Learning] almost everywhere.
    \item[Components of Machine Learning] \(\mathcal{A}\) takes \(\mathcal{D}\)
        and \(\mathcal{H}\) to get \(g\).
    \item[Machine Learning and Other Fields] related to DM, AI and statistics.
\end{description}

We have a frame, \((\mathcal{X}, \mathcal{Y}, f, \mathcal{D}, \mathcal{H},
\mathcal{A}, g)\), but not a single concrete learning model to put in it: we
never said what any \(h \in \mathcal{H}\) actually computes, nor how
\(\mathcal{A}\) searches. The next chapter fills in the simplest useful choice
of \(\mathcal{H}\) and \(\mathcal{A}\), for the simplest useful output space
\(\mathcal{Y} = \left\{ -1, +1 \right\}\) --- that is, learning to answer
yes/no.
